\chapter{Wiener identification: a nonlinear output map after linear dynamics}
\textit{Evidence basis: }\path|results/wiener_ramps_20260908_214558_458/REPORT.md|. The following detailed account is adapted from the archived study; references to simulation and verification describe historical work.\par

Fitted in MATLAB R2018b, System Identification Toolbox 9.9, using \path|nlhw|,
\path|unitgain| on the input and \path|poly1d(2)| on the output.

\section{Selected model and free-run fits}

{\scriptsize\setlength{\tabcolsep}{1.5mm}\begin{longtable}{>{\raggedright\arraybackslash}p{50.33mm}>{\raggedright\arraybackslash}p{50.33mm}>{\raggedright\arraybackslash}p{50.33mm}}\caption{Wiener identification: a nonlinear output map after linear dynamics: Selected model and free-run fits. Archived values.}\\\toprule
Record & Fit, NRMSE (\%) & RMSE (rpm)\\\midrule
\endfirsthead
Record & Fit, NRMSE (\%) & RMSE (rpm)\\\midrule
\endhead
Training: converged chirp prefix & 94.0903 & 21.4845\\
Ramp/hold model-selection record & 87.1724 & 39.5604\\
Held-out ramp/hold final validation & 87.0954 & 41.9286\\
\bottomrule\end{longtable}}

Selected orders \textbf{[nb nf nk] = [1 3 0]} at
Ts=0.015 s. nb is the number of numerator coefficients excluding delay
zeros; nf is denominator degree. Both are capped at 10 (the cap is per polynomial,
not a total of 10 across both). Output nonlinearity degree is exactly 2.
Maximum absolute linear-block pole is 0.9786791711; the selected
linear block is stable. Maximum absolute final-validation error is
58.5106 rpm.

Fit = 100 * (1 - norm(y-yhat)/norm(y-mean(y))). These are free-run simulation
scores with zero normalized initial states, not one-step prediction scores.
Measured validation speed is not fed back to the predictor or used to estimate
initial states. The coefficient-only implementation and MATLAB \path|compare| metric
were checked against the saved predictions; the exported metrics were independently
recomputed from the CSV files.

\section{Training and validation separation}

Training uses ONLY the first 1187 samples, 0-17.790 s, of the
previous 9000-10000 rpm chirp. Every sample at and after the first solver failure
at 17.805 s was excluded. Consequently the accepted chirp frequency coverage is
approximately 0.1-0.36685 Hz, not the requested full
0.1-1 Hz. The discarded samples were not interpolated, repaired or used for fitting.

Ramp/hold records are separate simulations, each with 60 s retained after a 30 s
preparation interval. The selection and final validation records use different
levels, ramp rates and hold durations. All ramp profiles stay within 9000-10000 rpm.
Across them, nonzero rate magnitudes are 50, 60, 95, 100 and 150 rpm/s.
Actual fuel flow Wf (lbm/s), rather than the speed-demand signal, is the model input;
mechanical shaft speed Nmech (rpm) is the output.

Only the chirp estimates model parameters and normalization. The selection ramp
chooses model order and whether training-only refinement is retained. The final
validation ramp is evaluated after that choice, without refitting.
44 candidate order/delay structures were tested, with up to
80 iterations per candidate and up to 200 for final refinement. The coarse search
tests all nb/nf pairs 1-4 and equal nb/nf orders 5-10, with delays 0 and 1; local
neighboring orders are then tested. "Best" means best selection fit among this
search, not a proof of a global optimum or an exhaustive search of all orders.
Candidate errors and scores are retained in candidate\_\allowbreak scores.csv.

\section{Accepted simulation checks}

{\scriptsize\setlength{\tabcolsep}{1.5mm}\begin{longtable}{>{\raggedright\arraybackslash}p{29.00mm}>{\raggedright\arraybackslash}p{29.00mm}>{\raggedright\arraybackslash}p{29.00mm}>{\raggedright\arraybackslash}p{29.00mm}>{\raggedright\arraybackslash}p{29.00mm}}\caption{Wiener identification: a nonlinear output map after linear dynamics: Accepted simulation checks. Archived values.}\\\toprule
Record & Max normalized flow residual & Max iterations & Min fuel (lbm/s) & Min stability margin (\%)\\\midrule
\endfirsthead
Record & Max normalized flow residual & Max iterations & Min fuel (lbm/s) & Min stability margin (\%)\\\midrule
\endhead
Selection & 9.9944e-10 & 30 & 0.9337 & 16.6316\\
Validation & 9.9968e-10 & 20 & 0.8548 & 15.9367\\
\bottomrule\end{longtable}}

Both entire 90 s source runs pass the 1e-9 flow tolerance and 200-iteration limit.
Recorded fuel remains positive and compressor stability margin remains positive.
The original GPT model and PI controller were not changed to obtain these runs.
These are noise-free synthetic closed-loop experiments. Fit on these trajectories
does not establish validity for failed high-frequency chirps, abrupt full-range
steps, untested operating conditions or measurement noise. The Wiener structure
is an approximation of the gas-turbine dynamics.

\section{Numerical model}

Let u\_\allowbreak n=(Wf-u0)/su; v=filter(B,F,u\_\allowbreak n); Nmech=y0+sy*polyval(p,v),
with zero filter initial states and coefficients ordered as MATLAB expects:

\begin{quote}\footnotesize
\begin{verbatim}
B = [1]
F = [1, -2.746098163319219, 2.550083469024761, -0.8028568299992156]
p = [-1.328480990017108e-07, 0.001843948442874731, -0.001341555354298684]  # [quadratic, 
linear, constant]
u0 = 0.9810869146816966 lbm/s
su = 1.108919409314424 lbm/s
y0 = 9000.000000000004 rpm
sy = 363.6986568026227 rpm
\end{verbatim}
\end{quote}

Leading zeros in B already implement nk. Do not add the delay again.
Normalization is part of the saved model wrapper; it must be applied exactly.
For other starting operating points, consistent filter initial states are needed.

\section{Use in MATLAB}

\begin{quote}\footnotesize
\begin{verbatim}
s = load('best_wiener_model.mat');
disp(s.bestModel)
% Full toolbox model, with the supplied normalized final validation data:
compare(s.zv, s.bestModel, Inf, compareOptions('InitialCondition','zero'))

% Coefficient-only implementation, physical fuel -> physical rpm:
d = readtable('validation_data.csv');
Nhat = predict_wiener_gpt(d.Wf, 'best_wiener_model.mat');
\end{verbatim}
\end{quote}

\section{Files}

\par\noindent\textbullet\ best\_\allowbreak wiener\_\allowbreak model.mat: selected idnlhw model, coefficients, scales, metrics,
iddata estimation/selection/validation objects and predicted outputs.
\par\noindent\textbullet\ model\_\allowbreak coefficients.json and fit\_\allowbreak metrics.json: portable numerical model and scores.
\par\noindent\textbullet\ identification\_\allowbreak data.mat: data tables, normalized iddata objects and scaling.
\par\noindent\textbullet\ training\_\allowbreak chirp\_\allowbreak valid.csv: only the accepted training segment.
\par\noindent\textbullet\ selection\_\allowbreak data.csv and validation\_\allowbreak data.csv: converged ramp/hold measurements.
\par\noindent\textbullet\ selection\_\allowbreak profile.csv and validation\_\allowbreak profile.csv: ramp breakpoints.
\par\noindent\textbullet\ *\_\allowbreak predictions.csv: actual output, simulated output and residuals.
\par\noindent\textbullet\ candidate\_\allowbreak scores.csv and search\_\allowbreak checkpoint.mat: search results and candidate models.
\par\noindent\textbullet\ wiener\_\allowbreak fit.png: training and held-out validation comparison and errors.
\par\noindent\textbullet\ experiments.mat: ramp data, metadata and original simulation-directory references.
\par\noindent\textbullet\ reproduction scripts and MATLAB logs: included for traceability. Simulation
scripts require the original project layout and adjacent T-MATS installation.
